\section{APPENDIX C}

\subsection{Paragraf 1: Eksperimen 2x2}
\begin{frame}{Prosedur Iteratif}
    \begin{itemize}
        \item \textbf{Tahap Awal:} Inisialisasi $\ket{+}$ dan proyeksi ke $\ket{11}$ untuk estimasi 2-bit.
        \item \textbf{Konvergensi:} Terminasi iterasi pada langkah ke-4 saat tercapai stabilitas vektor.
    \end{itemize}
\end{frame}

\begin{frame}{Pertanyaan yang dijawab pada Paragraf 1}
    \centering
    \small
    \begin{tabular}{lp{4.5cm}p{6cm}}
        \toprule
        Kalimat ke & Pertanyaan & Jawaban \\
        \midrule
        1 & Fokus bagian 1 sirkuit? & Estimasi 2-bit dari nilai eigen tertinggi. \\
        4 & Kapan iterasi berhenti? & Saat vektor hasil pengukuran sudah stabil (iterasi 4). \\
        \bottomrule
    \end{tabular}
\end{frame}

\subsection{Paragraf 2: Akurasi Fase}
\begin{frame}{Rotasi Pengukuran}
    \begin{itemize}
        \item \textbf{Basis x, y, r:} Rotasi basis pengukuran untuk meningkatkan resolusi data fase.
        \item \textbf{Tabel IV:} Outcome untuk setiap variasi rotasi basis.
    \end{itemize}
\end{frame}

\begin{frame}{Pertanyaan yang dijawab pada Paragraf 2}
    \centering
    \small
    \begin{tabular}{lp{4.5cm}p{6cm}}
        \toprule
        Kalimat ke & Pertanyaan & Jawaban \\
        \midrule
        1 & Akurasi ditingkatkan? & Via rotasi basis pengukuran sumbu x, y, dan r. \\
        2 & Sebaran data? & Rincian counts dan persentase pada Tabel IV. \\
        \bottomrule
    \end{tabular}
\end{frame}

\subsection{Paragraf 3: Hasil Akhir Eigen}
\begin{frame}{Estimasi Akhir}
    \begin{itemize}
        \item \textbf{Eigenvector:} Disajikan dalam persamaan (C1) dengan margin galat $\delta$.
        \item \textbf{Eigenvalue:} Nilai $\Lambda_{max} = 0.11$ dalam representasi biner.
    \end{itemize}
\end{frame}

\begin{frame}{Pertanyaan yang dijawab pada Paragraf 3}
    \centering
    \small
    \begin{tabular}{lp{4.5cm}p{6cm}}
        \toprule
        Kalimat ke & Pertanyaan & Jawaban \\
        \midrule
        1 & Hasil eigenvector? & Dinyatakan via Persamaan (C1) dengan galat $\delta$. \\
        2 & Nilai eigenvalue biner? & Estimasi $\Lambda_{max} = 0.11$. \\
        \bottomrule
    \end{tabular}
\end{frame}

\subsection{Paragraf 4: Eksperimen QPE 4-Qubit}
\begin{frame}{Fase Kedua Protokol}
    \begin{itemize}
        \item \textbf{Sistem 4-Qubit:} Register register eigenvalue 3-bit dan register eigenvector.
        \item \textbf{Kehilangan Koherensi:} Dekoherensi total pada hardware riil akibat kedalaman sirkuit.
        \item \textbf{Fidelitas Simulator:} Hasil simulator hampir ideal dengan fidelitas mencapai 0.965.
    \end{itemize}
\end{frame}

\begin{frame}{Pertanyaan yang dijawab pada Paragraf 4}
    \centering
    \small
    \begin{tabular}{lp{4.5cm}p{6cm}}
        \toprule
        Kalimat ke & Pertanyaan & Jawaban \\
        \midrule
        4 & Masalah utama QPE? & Kehilangan koherensi total pada hardware riil. \\
        5 & Kualitas simulator? & Hasil hampir ideal dengan fidelitas 0.965. \\
        6 & Estimasi QPE? & Biner 0.111 yang setara desimal 0.875. \\
        \bottomrule
    \end{tabular}
\end{frame}

\subsection{Paragraf 5: Skalabilitas 4x4}
\begin{frame}{Uji Sistem Multivariat}
    \begin{itemize}
        \item \textbf{Replikasi:} Replikasi prosedur untuk matriks 4x4 guna menguji skalabilitas.
        \item \textbf{Dekoherensi:} Dampak destruktif dekoherensi pada sistem dimensi tinggi.
        \item \textbf{Hasil 4x4:} Estimasi eigenvector utama dalam kombinasi linear basis.
    \end{itemize}
\end{frame}

\begin{frame}{Pertanyaan yang dijawab pada Paragraf 5}
    \centering
    \small
    \begin{tabular}{lp{4.5cm}p{6cm}}
        \toprule
        Kalimat ke & Pertanyaan & Jawaban \\
        \midrule
        1 & Pengujian 4x4? & Ikuti langkah kasus $2 \times 2$ secara bertahap. \\
        5 & Kondisi hardware? & Kedalaman sirkuit cegah perolehan hasil berguna. \\
        6 & Estimasi akhir? & Kombinasi linear basis $00, 01, 10, 11$. \\
        \bottomrule
    \end{tabular}
\end{frame}
